\documentclass{article}
\usepackage{import}
\import{../../template/}{article-template}
\graphicspath{{../../images}}
\addbibresource{../../template/library.bib}

\title{Partial Differential Equations}
\author{PENG}
\date{Created: 20260322, Version: 20260910}

\begin{document}
\maketitle
\tableofcontents

\section{Introduction}
\subsection{Notation}
\begin{definition}[$n$-dimensional real Euclidean space]
    $\mathbb{R}^{n}$
\end{definition}

\begin{definition}[point in $\mathbb{R}^{n}$]
    \begin{equation}
        \label{eq:x}
        x :=
        \begin{bmatrix}
            x_{1} \\ \vdots \\ x_{n}  \\
        \end{bmatrix}
    \end{equation}
\end{definition}

\begin{definition}[$u(x)$]
    \begin{equation}
        \label{eq:u}
        u : U \to \mathbb{R} \Leftrightarrow u(x) := u(x_{1}, \ldots, x_{n}) \quad (x \in U)
    \end{equation}
\end{definition}

\begin{definition}[standard coordinate vector]
    \begin{equation}
        \label{eq:ei}
        e_{i} :=
        \begin{bmatrix}
            0 \\ \vdots \\ 0 \\ 1 \\ \vdots \\ 0
        \end{bmatrix}
    \end{equation}
\end{definition}

\begin{definition}[derivative]
    \begin{equation}
        \label{eq:derivative}
        u_{x_{i}} := \frac{\partial u}{\partial x_{i}}(x) = \lim_{h \to 0} \frac{u(x+he_{i}) - u(x)}{h}
    \end{equation}
\end{definition}

\begin{definition}[smooth]
    $u$ is smooth provided $u$ is \textbf{infinitely differentiable}.
\end{definition}

\begin{definition}[multi-index of order]
    \begin{gather}
        \label{eq:alpha}
        \alpha = \begin{bmatrix} \alpha_{1} \\ \vdots \\ \alpha_{n} \end{bmatrix} \\
        \label{eq:|alpha|}
        |\alpha| = \alpha_{1} + \ldots + \alpha_{n}
    \end{gather}
    where $\alpha_{i} \geq 0$
\end{definition}

\begin{definition}[$D^{\alpha}$]
    \begin{equation}
        \label{eq:D-alpha}
        D^{\alpha} u(x) := \frac{\partial^{|\alpha|} u(x)}{\partial x_{1}^{\alpha_{1}} \ldots \partial x_{n}^{\alpha_{n}}}
        = \partial_{x_{1}}^{\alpha_{1}} \ldots \partial_{x_{n}}^{\alpha_{n}} u
    \end{equation}
\end{definition}

\begin{definition}[$D^{k}$]
    \begin{equation}
        \label{eq:Dk}
        D^{k} u(x) := \{D^{\alpha} u(x) : |\alpha| = k\}
    \end{equation}
\end{definition}

\begin{definition}[gradient vector]
    \begin{equation}
        \label{eq:D}
        Du = \nabla u :=
        \begin{bmatrix}
            \frac{\partial u}{\partial x_{1}} \\
            \vdots                            \\
            \frac{\partial u}{\partial x_{n}}
        \end{bmatrix}
        \in \mathbb{R}^{n}
    \end{equation}
\end{definition}

\begin{definition}[Hessian matrix]
    Real symmetrical matrices:
    \begin{equation}
        \label{eq:D2}
        D^{2}u :=
        \begin{bmatrix}
            u_{x_{1}x_{1}} & \dots  & u_{x_{1}x_{n}} \\
            \vdots         & \ddots & \vdots         \\
            u_{x_{n}x_{1}} & \dots  & u_{x_{n}x_{n}}
        \end{bmatrix}
    \end{equation}
\end{definition}

\begin{definition}[$A:B$]
    \begin{equation}
        \label{eq:A:B}
        A : B = \sum_{i=1}^{m}\sum_{j=1}^{n}a_{ij}b_{ij}
    \end{equation}
\end{definition}

\begin{definition}[$|A|$]
    \begin{equation}
        \label{eq:|A|}
        |A| = \sqrt{\sum_{i=1}^{m}\sum_{j=1}^{n} a_{ij}^{2}}
    \end{equation}
\end{definition}

\begin{definition}[vector dot product and matrix multiplication]
    \begin{equation}
        A \cdot B = A^{T} B =
        [a_{1} \dots a_{n}]
        \begin{bmatrix}
            b_{1} \\ \vdots \\ b_{n}
        \end{bmatrix}
        = \sum_{i=1}^{n} a_{i}b_{i}
    \end{equation}
\end{definition}

\begin{definition}[radial derivative]
    \begin{equation}
        \label{eq:radial-derivative}
        \begin{aligned}
            u_{r} := & \frac{x}{|x|} \cdot Du
            = \frac{1}{|x|} [x_{1} \dots x_{n}]
            \begin{bmatrix}
                \frac{\partial u}{\partial x_{1}} \\
                \vdots                            \\
                \frac{\partial u}{\partial x_{n}}
            \end{bmatrix}                                               \\
            =        & \frac{\sum_{i=1}^{n} x_{i}\frac{\partial u}{\partial x_{i}}}{\sqrt{\sum_{i=1}^{n} x_{i}^{2}}}
            \quad (x \neq 0)
        \end{aligned}
    \end{equation}
\end{definition}

\begin{definition}[$\boldsymbol{u}$]
    \begin{equation}
        \label{eq:us}
        \boldsymbol{u} : U \to \mathbb{R}^{m} \Leftrightarrow \boldsymbol{u} :=
        \begin{bmatrix}
            u^{1} \\ \vdots \\ u^{m}
        \end{bmatrix}
    \end{equation}
\end{definition}

\begin{definition}[$D^{\alpha}\boldsymbol{u}$]
    \begin{equation}
        \label{eq:Ds-alpha}
        D^{\alpha} \boldsymbol{u} :=
        \begin{bmatrix}
            D^{\alpha}u^{1} \\ \vdots \\ D^{\alpha}u^{n}
        \end{bmatrix}
    \end{equation}
\end{definition}

\begin{definition}[$D^{k}\boldsymbol{u}$]
    \begin{equation}
        \label{eq:Dsk}
        D^{k} \boldsymbol{u} := \{D^{\alpha} \boldsymbol{u} : |\alpha| = k\}
    \end{equation}
\end{definition}

\begin{definition}[gradient matrix]
    \begin{equation}
        \label{eq:Ds}
        D\boldsymbol{u} :=
        \begin{bmatrix}
            u_{x_{1}}^{1} & \dots  & u_{x_{n}}^{1} \\
            \vdots        & \ddots & \vdots        \\
            u_{x_{1}}^{m} & \dots  & u_{x_{n}}^{m}
        \end{bmatrix}
    \end{equation}
\end{definition}

\subsection{PDE}
\begin{definition}[k\textsuperscript{th}-order PDE]
    \begin{equation}
        \label{eq:pde}
        F(D^{k}u(x), D^{k-1}u(x), \dots, Du(x), u(x), x) = 0
    \end{equation}
    where $x \in U$, $u : U \to \mathbb{R}$, and $F : \mathbb{R}^{n^{k}} \times \mathbb{R}^{n^{k-1}} \times \dots \times \mathbb{R}^{n} \times \mathbb{R} \times U \to \mathbb{R}$
\end{definition}

\begin{figure}[h]
    \centering
    \includegraphics[width=6cm]{ns-equation.png}
    \caption{FINITE TIME BLOWUP FOR NAVIER–STOKES \cite{OPENAI_2026_FINITE_TIME_BLOWUP_FOR_NAVIERSTOKES}.}
    \label{fig:ns-equation}
\end{figure}

\begin{definition}[linear PDE]
    \begin{equation}
        \label{eq:linear-pde}
        \sum_{|\alpha| \leq k} a_{\alpha}(x) D^{\alpha} u = f(x)
    \end{equation}

    $f \equiv 0 \Rightarrow$ homogeneous
\end{definition}

\begin{definition}[semi-linear]
    \begin{equation}
        \label{eq:semi-linear}
        \sum_{|\alpha| = k} a_{\alpha}(x) D^{\alpha}u + a_{0} (D^{k-1}u, \dots, Du, u, x) = 0
    \end{equation}
\end{definition}

\begin{definition}[quasi-linear]
    \begin{equation}
        \label{eq:quasi-linear}
        \sum_{|\alpha| = k} a_{\alpha}(D^{k-1}u, \dots, Du, u, x) D^{\alpha}u + a_{0} (D^{k-1}u, \dots, Du, u, x) = 0
    \end{equation}
\end{definition}

\begin{remark}[nonlinear]
    Full non-linearity depends on the highest order derivatives.
\end{remark}

\begin{definition}[k\textsuperscript{th}-order system of PDE]
    \begin{equation}
        \label{eq:pde-system}
        F(D^{k}\boldsymbol{u}(x), D^{k-1}\boldsymbol{u}(x), \dots, D\boldsymbol{u}(x), \boldsymbol{u}(x), x) = 0
    \end{equation}
    where $x \in U$, $\boldsymbol{u} : U \to \mathbb{R}^{m}$, and $F : \mathbb{R}^{mn^{k}} \times \mathbb{R}^{mn^{k-1}} \times \dots \times \mathbb{R}^{mn} \times \mathbb{R}^{m} \times U \to \mathbb{R}^{m}$
\end{definition}

\subsection{Solutions to PDE}
\begin{definition}[solutions to PDE \cite{EVANS_2022_PARTIAL_DIFFERENTIAL_EQUATIONS}]~
    \begin{itemize}
        \item classic solutions
              \begin{itemize}
                  \item explicit, closed-form formula: polynomials, exp, cos, sin
                  \item or at least show such a solution exists, and to deduce various of its properties
              \end{itemize}
        \item weak solutions and regularity (smoothness)
    \end{itemize}
\end{definition}

\begin{remark}
    Existence is the prime attribute of what we study.
\end{remark}

\section{Four Fundamental Linear PDE}
\begin{definition}[Fundamental linear PDE]~
    \begin{itemize}
        \item Transport equation
        \item Laplace's equation
        \item Heat equation
        \item Wave equation
    \end{itemize}
\end{definition}

\begin{lemma}
    $\frac{dy}{d(x+b)} = \frac{dy}{dx}$
\end{lemma}

\begin{proof}
    \begin{equation*}
        \frac{dy}{d(x+b)} = \frac{\frac{dy}{dx}}{\frac{d(x+b)}{dx}} = \frac{dy}{dx}
    \end{equation*}
\end{proof}

\subsection{Transport Equation}
\begin{definition}[Transport equation]
    \begin{equation}
        \label{equ:transport-equation}
        u_{t} + b \cdot Du =
        \frac{\partial u}{\partial t} +
        [b_{1} b_{2} \dots b_{n}]
        \begin{bmatrix}
            \frac{\partial u}{\partial x_{1}} \\
            \frac{\partial u}{\partial x_{2}} \\
            \vdots                            \\
            \frac{\partial u}{\partial x_{n}}
        \end{bmatrix}
        = \frac{\partial u}{\partial t} + \sum_{i=n}^{n}b_{i}\frac{\partial u}{\partial x_{i}} = 0
    \end{equation}
\end{definition}

\begin{example}
    For 1D and 2D cases:
    \begin{gather*}
        \frac{\partial u}{\partial t} + b \frac{\partial u}{\partial x} = 0 \\
        \frac{\partial u}{\partial t} + b_{1} \frac{\partial u}{\partial x} + b_{2} \frac{\partial u}{\partial y} = 0
    \end{gather*}
\end{example}

\begin{theorem}
    \label{theo:transport-equation-solution}
    For Eq. \ref{equ:transport-equation}, $u$ is constant on the line $(x,t) + (b,1)s$ for $s \in \mathbb{R}$.
\end{theorem}

\begin{proof}
    Let $z(s) := u(x+bs,t+s)$.
    \begin{gather*}
        \frac{dz}{ds} = \frac{\partial u}{\partial (t+s)}\frac{d(t+s)}{ds} + \frac{\partial u}{\partial (x+sb)}\frac{d (x+sb)}{ds}
        = \frac{\partial u}{\partial t} + b\frac{\partial u}{\partial x}
        \overset{\text{Eq. \ref{equ:transport-equation}}}{=} 0 \\
        z(s) = u(x+sb,t+s) = C
    \end{gather*}
\end{proof}

\begin{proposition}[initial-value problem]
    \begin{equation*}
        \begin{cases}
            \frac{\partial u}{\partial t} + b\frac{\partial u}{\partial x} = 0 \\
            u(x,0) = g(x)
        \end{cases}
    \end{equation*}

    \begin{equation*}
        u(x,t) \overset{\text{Theo. \ref{theo:transport-equation-solution}}}{=} u(x-tb,0) = g(x-tb)
    \end{equation*}
\end{proposition}

\begin{proposition}[nonhomogenous problem]
    \begin{equation*}
        \begin{cases}
            \frac{\partial u}{\partial t} + b \cdot \frac{\partial u}{\partial x} = f(x,t) \\
            u(x,0) = g(x)
        \end{cases}
    \end{equation*}

    Let $z(s) := u(x+sb,t+s)$.

    \begin{gather*}
        \frac{dz}{ds} = f(x+sb,t+s) \\
        z(s) = u(x+sb,t+s) = \int f(x+sb,t+s) ds \\
        z(0) - z(-t) = u(x,t) - u(x-tb,0) = u(x,t) - g(x-tb) = \int_{-t}^{0} f(x+sb,t+s) ds \\
        u(x,t) = g(x-tb) + \int_{-t}^{0} f(x+sb,t+s) ds
    \end{gather*}
\end{proposition}

\subsection{Laplace's Equation}
\begin{definition}[Laplace's equation]
    \begin{equation}
        \label{eq:laplace-equation}
        \Delta u = \nabla \cdot \nabla u = \nabla^{2} u =
        \left[ \frac{\partial}{\partial x_{1}} \frac{\partial}{\partial x_{2}} \dots \frac{\partial}{\partial x_{n}} \right]
        \begin{bmatrix}
            \frac{\partial u}{\partial x_{1}} \\
            \frac{\partial u}{\partial x_{2}} \\
            \vdots                            \\
            \frac{\partial u}{\partial x_{n}}
        \end{bmatrix}
        = \sum_{i=1}^{n} \frac{\partial^{2} u}{\partial x_{i}^{2}} = 0
    \end{equation}
\end{definition}

\begin{example}
    For 1D and 2D cases:
    \begin{equation*}
        \frac{\partial^2 u}{\partial x^{2}} = 0
    \end{equation*}

    \begin{equation*}
        \frac{\partial^2 u}{\partial x^{2}} + \frac{\partial^2 u}{\partial y^{2}} = 0
    \end{equation*}
\end{example}

\begin{definition}[outward derivative]
    \begin{equation*}
        \frac{\partial u}{\partial v} := \boldsymbol{v} \cdot Du =
        [v_{1} v_{2} \dots v_{n}]
        \begin{bmatrix}
            \frac{\partial u}{\partial x_{1}} \\
            \frac{\partial u}{\partial x_{2}} \\
            \vdots                            \\
            \frac{\partial u}{\partial x_{n}} \\
        \end{bmatrix}
        = \sum_{i=1}^{n} v_{i} \frac{\partial u}{\partial x_{i}}
    \end{equation*}
    where $\boldsymbol{v}$ is the unit outer normal vector of $\partial U$.
\end{definition}

\begin{figure}[h]
    \centering
    \includegraphics{outward-derivative.png}
    \caption{Outward derivative.}
    \label{fig:outward-derivative}
\end{figure}

\begin{example}
    For 1D and 2D cases:
    \begin{equation*}
        \frac{\partial u}{\partial v} = v\frac{\partial u}{\partial x}
    \end{equation*}

    \begin{equation*}
        \frac{\partial u}{\partial v} = v_{1}\frac{\partial u}{\partial x} + v_{2}\frac{\partial u}{\partial y}
    \end{equation*}
\end{example}

\begin{theorem}[Gauss-Green theorem]
    \begin{equation}
        \begin{aligned}
              & \int_{U} \nabla \cdot \boldsymbol{u} dx =
            \int_{U}
            \left[ \frac{\partial}{\partial x_{1}} \frac{\partial}{\partial x_{2}} \dots \frac{\partial}{\partial x_{n}} \right]
            \begin{bmatrix}
                u_{1}  \\
                u_{2}  \\
                \vdots \\
                u_{n}
            \end{bmatrix} dx
            = \int_{U} \sum_{i=1}^{n} \frac{\partial u_{i}}{\partial x_{i}} dx \\
            = & \int_{\partial U} \boldsymbol{u} \cdot \boldsymbol{v} dS
            = \int_{\partial U}
            \left[ u_{1} u_{2} \dots u_{3} \right]
            \begin{bmatrix}
                v_{1}  \\
                v_{2}  \\
                \vdots \\
                v_{n}
            \end{bmatrix} dS
            = \int_{\partial U} \sum_{i=1}^{n} u_{i}v_{i} dS
        \end{aligned}
    \end{equation}
\end{theorem}

\begin{example}
    For 1D and 2D cases:
    \begin{gather*}
        \int_{U} \frac{\partial u}{\partial x} dx = \int_{\partial U} uv dS \\
        \int_{U} \frac{\partial u_{1}}{\partial x} + \frac{\partial u_{2}}{\partial y} dx =
        \int_{\partial U} u_{1}v_{1} + u_{2}v_{2} dS
    \end{gather*}
\end{example}

Fundamental solution:
\begin{gather*}
    u(x) = v(r) \\
    r = |x| = \sqrt{\sum_{i=1}^{n}x_{i}^{2}} \\
    \frac{\partial r}{\partial x_{i}} = \frac{x_{i}}{r} \\
    \frac{\partial u}{\partial x_{i}} = \frac{dv}{dr} \frac{x_{i}}{r} \\
    \frac{\partial^{2} u}{\partial x_{i}^{2}} = \frac{d^{2}v}{dr^{2}} \frac{x_{i}^{2}}{r^{2}}
    + \frac{dv}{dr} \frac{r^{2} - x_{i}^{2}}{r^{3}} \\
    \Delta u = \sum_{i=1}^{n} \frac{\partial^{2} u}{\partial x_{i}^{2}} = \frac{d^{2}v}{dr^{2}} + \frac{dv}{dr}\frac{n-1}{r} = 0
\end{gather*}

For $n = 1$:

\begin{equation*}
    \frac{d^{2}v}{dr^{2}} = 0
\end{equation*}

\begin{equation*}
    v(r) = r + C
\end{equation*}

For $n \geq 2$:

\begin{equation*}
    \frac{d^{2}v}{dr^{2}} + \frac{dv}{dr}\frac{n-1}{r} = 0
\end{equation*}

\begin{equation*}
    \frac{v''(r)}{v'(r)} = \ln\left(\left| v'(r) \right|\right)' = \frac{1-n}{r}
\end{equation*}

\begin{equation*}
    \ln\left(\left| v'(r) \right|\right) = (1-n)\ln(r) + \ln(C) = \ln\left( \frac{C}{r^{n-1}} \right)
\end{equation*}

\begin{equation*}
    \left| v'(r) \right| = \frac{C}{r^{n-1}}
\end{equation*}

\begin{equation*}
    v'(r) = \frac{C}{r^{n-1}}
\end{equation*}

For $n = 2$:

\begin{equation*}
    v'(r) = \frac{C}{r}
\end{equation*}

\begin{equation*}
    v(r) = C_{1}\ln(r) + C_{2}
\end{equation*}

For $n \geq 3$:

\begin{equation*}
    v(r) = \frac{C_{1}}{2-n} \frac{1}{r^{n-2}} + C_{2} = \frac{C_{1}}{r^{n-2}} + C_{2}
\end{equation*}

\begin{definition}[fundamental solution of Laplace's equation]
    \begin{equation}
        \label{eq:laplace-equation-solution}
        \Phi(|x|) :=
        \begin{cases}
            -\frac{1}{2\pi}\ln|x| \quad                        & (n = 2)    \\
            \frac{1}{n(n-2)\alpha(n)}\frac{1}{|x|^{n-2}} \quad & (n \geq 3)
        \end{cases}
    \end{equation}
\end{definition}

\begin{definition}[Poisson's equation]
    \begin{equation}
        \label{eq:poisson-equation}
        -\Delta u = -\sum_{i=1}^{n} \frac{\partial^{2} u}{\partial x_{i}^{2}} = f
    \end{equation}
\end{definition}

\begin{theorem}
    Solution of Poisson's equation:
    \begin{equation}
        \label{eq:poisson-equation-solution}
        u(x) = \int_{y \in \mathbb{R}^{n}} \Phi(|x-y|) f(y) dy =
        \begin{cases}
            -\frac{1}{2\pi} \int_{\mathbb{R}^{n}} \ln\left( |x-y| \right) f(y) dy \quad       & (n = 2)    \\
            \frac{1}{n(n-2)\alpha(n)} \int_{\mathbb{R}^{n}} \frac{f(y)}{|x-y|^{n-2}} dy \quad & (n \geq 3)
        \end{cases}
    \end{equation}
\end{theorem}

\section{Asymptotics}
\subsection{Singular Perturbations}

\printbibliography
\end{document}
